\section{Implementação na placa}
\label{sec:placa}

\subsection{Pinagem}

A Tabela~\ref{tab:pinos} lista as restrições do arquivo
\sig{constraints/cpu\_top.ucf}, conforme o manual da Spartan-3E Starter Kit.

\begin{table}[H]
\centering
\caption{Pinagem utilizada (Spartan-3E Starter Kit).}
\label{tab:pinos}
\small
\begin{tabularx}{\linewidth}{>{\raggedright}p{3.0cm} >{\raggedright}p{3.1cm} l >{\raggedright\arraybackslash}X}
\toprule
Sinal & Pino & Padrão de I/O & Função \\
\midrule
\sig{CLK\_50MHZ} & C9  & LVCMOS33 & clock de \SI{50}{\mega\hertz} \\
\sig{SW<0>}      & L13 & LVTTL    & chave SW0: reset da CPU \\
\sig{LED<7>} \ldots \sig{LED<0>} & F9, E9, D11, C11, F11, E11, E12, F12 & LVTTL & LEDs LD7 a LD0 \\
\sig{LCD\_E}     & M18 & LVCMOS33 & \emph{enable} do LCD \\
\sig{LCD\_RS}    & L18 & LVCMOS33 & comando (0) / dado (1) \\
\sig{LCD\_RW}    & L17 & LVCMOS33 & sempre 0 (escrita) \\
\sig{SF\_D<8>} \ldots \sig{SF\_D<11>} & R15, R16, P17, M15 & LVCMOS33 & dados do LCD (4 bits) \\
\sig{SF\_CE0}    & D16 & LVCMOS33 & desabilita a StrataFlash (1) \\
\bottomrule
\end{tabularx}
\end{table}

\subsection{Procedimento}

\begin{enumerate}[nosep]
  \item Criar o projeto no ISE para o dispositivo XC3S500E, encapsulamento
        FG320, \emph{speed grade} $-4$, adicionando os arquivos de \sig{src/} e
        o \sig{.ucf}, com \sig{cpu\_top} como topo.
  \item Sintetizar, implementar e gerar o \emph{bitstream}; gravar a FPGA pelo
        iMPACT via USB.
  \item Com a chave SW0 em 1, a CPU fica parada no estado inicial e o LCD
        mostra \texttt{NOP} e \texttt{MEM[255] = 000}. Ao levar SW0 para 0, o
        programa começa a executar, uma instrução a cada \SI{2}{\second}.
  \item Para trocar o programa, substitui-se o bloco de inicialização em
        \sig{RAM.vhd} e repete-se a síntese. Para reiniciar o mesmo programa,
        basta voltar SW0 a 1 e depois a 0 --- lembrando que, como a RAM não é
        reiniciada pelo reset, dados que o programa tenha escrito na memória
        permanecem até que a FPGA seja regravada.
\end{enumerate}

\subsection{Utilização de recursos}

\pendente{Inserir aqui o resumo do relatório de síntese do ISE
(\emph{Design Summary}): número de \emph{slices}, flip-flops, LUTs, block RAMs
e a frequência máxima estimada.}

\subsection{Registro da execução}

\pendente{Inserir fotos da placa executando o Fibonacci (por exemplo, no
início, no meio da sequência e no \texttt{HALT} com \texttt{MEM[255] = 233}) e o
link do vídeo no YouTube (não listado).}

% Exemplo de inclusao de foto:
% \begin{figure}[H]
%   \centering
%   \includegraphics[width=0.7\linewidth]{fotos/fib_halt.jpg}
%   \caption{Placa ao final do Fibonacci: LCD mostrando HALT e MEM[255] = 233.}
% \end{figure}
% Video: \url{https://youtu.be/...}
